\documentclass[11pt]{article}
\usepackage[T1]{fontenc}
\usepackage[utf8]{inputenc}
\usepackage[margin=1in]{geometry}
\usepackage{CJKutf8}

\title{Accents, Ligatures and Non-Latin Text}
\author{Zoë Brontë, José Núñez, and Søren Kierkegård}
\date{12 septembre 2026}

\begin{document}
\maketitle

\section{Accented Latin text (UTF-8 source)}
Café, naïve, résumé, façade, crème brûlée, jalapeño, piñata, señor,
Ångström, smörgåsbord, Ærø, Œuvre, Straße, Zürich, Košice, Łódź,
Gdańsk, Ørsted, Þórshöfn, Ísafjörður.

The same words through \LaTeX{} accent commands:
Caf\'e, na\"\i ve, r\'esum\'e, fa\c{c}ade, cr\`eme br\^ul\'ee, jalape\~no,
\AA ngstr\"om, sm\"org\aa sbord, \AE r\o, \OE uvre, Stra\ss e, Z\"urich,
Ko\v{s}ice, \L\'od\'z, Gda\'nsk, \O rsted, \TH \'orsh\"ofn.

\section{Ligatures, dashes and quotes}
Five fluffy waffles offer efficient office affluence --- ff, fi, fl, ffi, ffl.
En dash: pages 12--15; em dash --- like this; ``double quotes'' and `single
quotes'; ellipsis\dots; 25\,\% and \S\,4; \copyright~2026; \pounds 75.

\section{Greek and Cyrillic words in text}
The Greek word for typesetting is stoicheiothesía, and the Cyrillic spelling
of Moscow is Moskva. Their native forms need a font with those scripts and
are therefore given only in transliteration in this T1-encoded document.

\section{CJK via inputenc utf8}
\begin{CJK}{UTF8}{min}
東京は日本の首都です。
\end{CJK}
(``Tokyo is the capital of Japan.'')

\begin{CJK}{UTF8}{gbsn}
排版是一门艺术。
\end{CJK}
(``Typesetting is an art.'')

\section{Mixed script paragraph}
The r\'esum\'e of Søren --- filed under ``K'' --- lists Zürich, Łódź and
\begin{CJK}{UTF8}{min}東京\end{CJK} as places of residence, 2019--2026.

\end{document}
